% Elaborado por Fabian Gomez
% Version 1.0

\section{References}

Esta sección lista los documentos normativos y guías técnicas que enmarcan la estructura, el contenido mínimo y las buenas prácticas aplicadas en este SyRS/SRS del Rover Olympus. \cite{ieee_29148, segoldmine_ppi} 
Los borradores (main/main1/main3) ya referencian ISO/IEC/IEEE 29148 y el NASA Systems Engineering Handbook como marco metodológico general, y el uso de CSP (CubeSat Space Protocol) como referencia técnica para el encapsulamiento de telemetría/comandos. \cite{ieee_29148, nasa_se_handbook} 
Adicionalmente, debido a que el SRS final incorpora un ajuste realista del UC‑05 hacia ``Safe Deactivation / Passivation'' y menciona consideraciones de protección planetaria (no interferencia con futuras misiones), se incluyen referencias NASA de Planetary Protection como documentos externos aplicables/consultados. \cite{nasa_se_handbook, ieee_29148, segoldmine_ppi, nasa_planetary_protection}

\subsection{Applicable Standards}

Los siguientes estándares y guías se consideran aplicables o de referencia directa para la elaboración y mantenimiento del SyRS/SRS:

\begin{enumerate}
    \item \textbf{ISO/IEC/IEEE 29148:2018 – Systems and software engineering — Life cycle processes — Requirements engineering.} \cite{ieee_29148}
    \begin{itemize}
        \item Aplica como marco para estructura del documento, criterios de ``buen requisito'', contenido esperado y prácticas de gestión/ingeniería de requisitos.
    \end{itemize}

    \item \textbf{NASA Systems Engineering Handbook – NASA SP‑2016‑6105 Rev2.} \cite{nasa_se_handbook}
    \begin{itemize}
        \item Aplica como guía de buenas prácticas SE (incluyendo relación requisitos–arquitectura–V\&V, y recomendaciones de tailoring para el contexto de proyecto).
    \end{itemize}

    \item \textbf{CubeSat Space Protocol (CSP) Specification.} \cite{csp_spec}
    \begin{itemize}
        \item Aplica como referencia técnica para el encapsulamiento y verificación de integridad de paquetes de telemetría/comandos en la interfaz rover–estación terrena, tal como se utiliza en el stack descrito en los borradores. \cite{ieee_29148, nasa_se_handbook, segoldmine_ppi}
    \end{itemize}
\end{enumerate}
